%!%
% Copyright (c) 2026 mingcheng <mingcheng@outlook.com>
% 
% File: preface.tex
% Author: mingcheng <mingcheng@outlook.com>
% File Created: 2026-03-29 00:58:27
% 
% Modified By: mingcheng <mingcheng@outlook.com>
% Last Modified: 2026-03-29 07:14:47
%%

% !TeX root = ../prompt-engineering-guide-for-beginners.tex

\chapter*{前言}

\section*{关于作者}

大家好，我是吕峰军（明城），一个写了十几年代码的技术人\footnote{\href{https://github.com/mingcheng}{我的 GitHub 主页}}。
平时在互联网公司做后端和架构相关的工作，同时也是一个多年的开源爱好者。
这几年 AI 的发展太快了，我也一直在关注和尝试各种 AI 工具。从最初的新奇到日常的依赖，过程中踩了不少坑，也攒下了一些经验。

\section*{为什么写这本小册子？}

写这本小册子的起因特别朴素。

公司里越来越多的小伙伴开始用 AI 了，但大家的使用效果参差不齐：有人用得飞起，三分钟搞定一篇周报；有人折腾半天，AI 给出来的东西完全不能用。

问题出在哪儿？大多数时候，不是 AI 不行，而是我们跟它「说话的方式」不对。
同样一个需求，换一种问法，结果可能天差地别。
这就是提示工程（Prompt Engineering）的价值所在：学会跟 AI 好好说话，让它真正成为你的得力助手。

这本小册子最初源自团队内部的分享。后来有同事说，身边很多朋友都有同样的困惑，建议我整理得更详细些，分享给更多人。
于是我趁着业余时间不断补充和完善，最终形成了你现在看到的这本小册子。

不管你是刚接触 AI 的纯新手，还是已经用了一阵但总觉得「差点意思」的老用户，我都希望这本小册子能帮到你。

\section*{怎么读这本小册子？}

「提示工程」听起来有点高大上，但它的底层逻辑并不玄乎。本质上它是用自然语言作为代码，对概率庞大的神经网络模型进行一次“黑盒编程”。
这不是什么需要门槛的前沿技术，而是一项极为实用的逻辑外脑驾驭技能。你不需要熟练掌握任何编程语言，不需要弄懂 Transformer 的底层结构，只要会打字、拥有清晰的结构化思维就能学会。

当然，掌握它也需要一些时间和逻辑训练。刚开始可能会觉得有点抽象甚至是“玄学”，但别担心，书里的内容摒弃了艰涩的算法推导，尽量用通俗易懂的工程化类比，并配合了大量的真实示例和可复用的模板，帮你快速上手这门“赛博手艺”。

在开始之前，我想给你三个建议：

\textbf{第一，别光看，要动手。}每一章里的例子和模板，都建议你打开 AI 工具实际试一遍。看十遍不如自己写一遍，亲手体验一次胜过读十页理论。

\textbf{第二，别追求完美，先用起来。}提示词没有标准答案，同一个需求可以有很多种写法。先从简单的开始，慢慢迭代优化，你会越用越顺手。

\textbf{第三，保持好奇心，但不必焦虑。}AI 模型和提示技巧都在飞速演进，你不需要追逐每一个新概念。只要掌握了核心原则，新工具出来时你也能很快上手。

\section*{关于更新与贡献}

AI 领域变化太快了，所以这本小册子也不会一成不变。我会持续更新内容，跟上最新的模型能力和技术发展。

这本小册子在 GitHub 上完全开源（基于 Apache 2.0 许可证）。
如果你发现了错误、有更好的案例、或者想补充新的内容，非常欢迎通过 Pull Request 来贡献。
同样，如果你有任何意见和建议，也欢迎提 Issue 告诉我。

众人拾柴火焰高，期待和你一起把这份资料做得更好。

\section*{致谢}

感谢公司同事们在日常工作中的实践和反馈，书中很多真实案例都来自大家的使用经验。
也感谢每一位愿意花时间阅读这本小册子的你。

希望它能帮你在 AI 时代多一项趁手的技能。那我们开始吧！

\section*{关键词}

提示工程（Prompt Engineering）、提示词（Prompt）、大语言模型（LLM）、人工智能（AI）、ChatGPT、零样本提示（Zero-shot）、少样本提示（Few-shot）、思维链（Chain of Thought）、多模态（Multimodal）、检索增强生成（RAG）
